% !TEX program = lualatex
% ---------------------------------------------------------------------------------------------
%  Readback note 19: Primitive labels of general linear groups over finite fields: the share of bad types and the cubic twist
%
%  Build:  python3 tools/rb.py note <slug>
%
%  One unit of a Lean artifact, rewritten as ordinary mathematics. The prose is hand-written. Lean names,
%  line ranges and counts come from data.tex; quoted Lean statements are copied from the artifact by the build.
%
%  Rules of the text: standard terminology; every symbol is explained where it first appears and listed again
%  in the notation at the end; no letter has two meanings; a proof is complete or says what it takes for granted;
%  Lean names appear only in the LEAN lines and in the concordance.
%
%  The signature rule (the build enforces it): a function, relation or predicate that the artifact defines is
%  written only through its notation macro, which the build makes from the spec (`notation`) with one slot for
%  each explicit argument of the Lean definition. No argument is fixed and dropped, also not "for brevity" and
%  not for one section. Every argument and every hypothesis of a stated declaration stands inside its
%  statement, or in a Setting that the statement cites by number; only implicit and instance arguments (the
%  ambient objects) may stand in the conventions.
% ---------------------------------------------------------------------------------------------
\documentclass[11pt,a4paper]{article}
\ifdefined\pdfvariable\pdfvariable suppressoptionalinfo \numexpr 1+2+4+8\relax\fi
\usepackage[a4paper,hmargin=35mm,top=22mm,bottom=27mm,footskip=12mm,nohead]{geometry}

\usepackage{fontspec}
\directlua{
  luaotfload.add_fallback("gfallback", {"STIXTwoMath-Regular.otf:mode=node;", "DejaVuSans.ttf:mode=node;"})
  luaotfload.add_fallback("gmonofallback", {"DejaVuSansMono.ttf:mode=node;", "STIXTwoMath-Regular.otf:mode=node;"})
}
\setmainfont{NewCM10}[Extension=.otf,UprightFont=*-Book,ItalicFont=*-BookItalic,
  BoldFont=*-Bold,BoldItalicFont=*-BoldItalic,SmallCapsFont=*-Book,SmallCapsFeatures={Letters=SmallCaps},
  RawFeature={fallback=gfallback}]
\setsansfont{NewCMSans10}[Extension=.otf,UprightFont=*-Book,ItalicFont=*-BookOblique,
  BoldFont=*-Bold,BoldItalicFont=*-BoldOblique,RawFeature={fallback=gfallback}]
\setmonofont{JetBrainsMono}[Extension=.otf,UprightFont=*-Regular,ItalicFont=*-Italic,
  BoldFont=*-SemiBold,BoldItalicFont=*-SemiBoldItalic,Scale=MatchLowercase,
  RawFeature={-calt;-liga;fallback=gmonofallback}]
\usepackage{amsmath}
\usepackage{amsthm}
\usepackage[math-style=TeX]{unicode-math}
\setmathfont{NewCMMath-Book.otf}
\emergencystretch=2.5em\relax

\usepackage{xcolor}
\usepackage{tikz}
\usetikzlibrary{calc,arrows.meta,positioning,backgrounds}
\usepackage[protrusion=true,expansion=true]{microtype}
\usepackage{enumitem}
\usepackage{needspace}
\usepackage{booktabs}
\usepackage{array}
\usepackage{longtable}
\usepackage{fancyhdr}
\usepackage{lastpage}
\usepackage{multicol}
\usepackage[hidelinks,pdfusetitle]{hyperref}

\definecolor{ink}{HTML}{0F172A}
\definecolor{inksoft}{HTML}{475569}
\definecolor{inkfaint}{HTML}{94A3B8}
\definecolor{rulegrey}{HTML}{CBD5E1}
\definecolor{arith}{HTML}{C2410C}
\definecolor{est}{HTML}{7C3AED}
\definecolor{ana}{HTML}{0891B2}
\definecolor{ok}{HTML}{047857}
\definecolor{leancol}{HTML}{3730A3}

\title{Primitive labels of general linear groups over finite fields: the share of bad types and the cubic twist}
\author{grh-compact}
\date{}

% ------------------------------------------------------------------------------------- macros
\newcommand{\leanfont}{\ttfamily\fontsize{8.6}{10.5}\selectfont\color{leancol}}
\newcommand{\lean}[1]{\csname LN@#1\endcsname}
\newcommand{\fillin}[1]{{\color{arith}\itshape[#1]}}  % a place of the template that is not written yet; the build reports every one that is left
\newcommand{\leanquote}[1]{\csname LQ@#1\endcsname}
% generated by rb.py note; do not edit
\newcommand{\ConcDecls}{15}
\newcommand{\Pin}{72fead10}
\newcommand{\PinText}{that of commit {\leanfont 72fead10}}
\newcommand{\ArtifactName}{grh-compact}
\expandafter\newcommand\csname LN@IsPrimitiveLabel\endcsname{\href{https://github.com/lsernau/grh-lean/blob/72fead102f663d5c7a407964a13dc0cf9d977207/grh-compact/src/New/Trace/BadShare.lean\#L29-L30}{{\leanfont IsPrimitiveLabel}}}
\expandafter\newcommand\csname LN@scaleRoots\endcsname{\href{https://github.com/lsernau/grh-lean/blob/72fead102f663d5c7a407964a13dc0cf9d977207/grh-compact/src/New/Trace/BadShare.lean\#L35-L41}{{\leanfont IsPrimitiveLabel.\allowbreak{}scaleRoots}}}
\expandafter\newcommand\csname LN@isPrimitiveRootprodroots\endcsname{\href{https://github.com/lsernau/grh-lean/blob/72fead102f663d5c7a407964a13dc0cf9d977207/grh-compact/src/New/Trace/BadShare.lean\#L45-L51}{{\leanfont IsPrimitiveLabel.\allowbreak{}isPrimitiveRoot\_prod\_roots}}}
\expandafter\newcommand\csname LN@goodTuples\endcsname{\href{https://github.com/lsernau/grh-lean/blob/72fead102f663d5c7a407964a13dc0cf9d977207/grh-compact/src/New/Trace/BadShare.lean\#L55-L57}{{\leanfont goodTuples}}}
\expandafter\newcommand\csname LN@existslabelIdemeqoffirreducible\endcsname{\href{https://github.com/lsernau/grh-lean/blob/72fead102f663d5c7a407964a13dc0cf9d977207/grh-compact/src/New/Trace/BadShare.lean\#L64-L86}{{\leanfont exists\_labelIdem\_eq\_off\_irreducible}}}
\expandafter\newcommand\csname LQ@existslabelIdemeqoffirreducible\endcsname{\leanquoteline{0}{lemma exists\_labelIdem\_eq\_off\_irreducible \{k n a\} [Field k] [Finite k] [Fintype n] (hr :}\leanquoteline{2}{Fintype.card n = 3 \textasciicircum{} (a + 1)) (hq : 9 ∣ .card k - 1) c \{S T : Set k[X]\} (hS : S.Subsingleton)}\leanquoteline{2}{(hT : T.Subsingleton) : ∃ P, (IsPrimitiveLabel (Fintype.card n) c →}\leanquoteline{2}{IsPrimitiveLabel (Fintype.card n) P ∧ P ∉ S ∪ T) ∧ ∀ g, ¬Irreducible g.1.charpoly →}\leanquoteline{2}{(labelIdem n c).coeff g = (labelIdem n P).coeff g}}
\expandafter\newcommand\csname LN@badsharele\endcsname{\href{https://github.com/lsernau/grh-lean/blob/72fead102f663d5c7a407964a13dc0cf9d977207/grh-compact/src/New/Trace/BadShare.lean\#L92-L136}{{\leanfont bad\_share\_le}}}
\expandafter\newcommand\csname LQ@badsharele\endcsname{\leanquoteline{0}{lemma bad\_share\_le \{n ι\} [Fintype n] [Fintype ι] \{k : ι → \_\}}\leanquoteline{4}{[∀ w, Field (k w)] [∀ w, Finite (k w)] \{q M r f : ℕ\} [NeZero M] (ε L : ℝ) (hq : ∀ w, .card (k w)}\leanquoteline{4}{= q) (hr : card n = r) (hf : card ι = f) (hn : 1 < r) (h2 : r < q) (hε : ε ≤ 1)}\leanquoteline{4}{(hs : (1 + M + 2 * r) / q ≤ ε) (hL : 2 * r * (2 * (r * .log q) + 1) * L ≤ f) :}\leanquoteline{4}{((1 - ∑ c ∈ goodTuples M r, tupleIdem (k := k) n c).coeff 1).re ≤ f * ε + exp (-L)}}
\expandafter\newcommand\csname LN@IsLabel\endcsname{\href{https://github.com/lsernau/grh-lean/blob/72fead102f663d5c7a407964a13dc0cf9d977207/grh-compact/src/Lib/LocalGL/Green/FrobeniusOrbits.lean\#L48-L48}{{\leanfont IsLabel}}}
\expandafter\newcommand\csname LN@labelIdem\endcsname{\href{https://github.com/lsernau/grh-lean/blob/72fead102f663d5c7a407964a13dc0cf9d977207/grh-compact/src/Lib/LocalGL/Green/FrobeniusOrbits.lean\#L61-L63}{{\leanfont labelIdem}}}
\expandafter\newcommand\csname LN@tupleIdem\endcsname{\href{https://github.com/lsernau/grh-lean/blob/72fead102f663d5c7a407964a13dc0cf9d977207/grh-compact/src/Lib/LocalGL/Green/LabelIdempotents.lean\#L108-L108}{{\leanfont tupleIdem}}}
\expandafter\newcommand\csname LN@isCentralProjsumtupleIdem\endcsname{\href{https://github.com/lsernau/grh-lean/blob/72fead102f663d5c7a407964a13dc0cf9d977207/grh-compact/src/Lib/LocalGL/Green/LabelIdempotents.lean\#L124-L125}{{\leanfont isCentralProj\_sum\_tupleIdem}}}
\expandafter\newcommand\csname LN@labelChardegreebounds\endcsname{\href{https://github.com/lsernau/grh-lean/blob/72fead102f663d5c7a407964a13dc0cf9d977207/grh-compact/src/Lib/LocalGL/Green/LabelDegree.lean\#L35-L57}{{\leanfont labelChar\_degree\_bounds}}}
\expandafter\newcommand\csname LN@labelCharmuldet\endcsname{\href{https://github.com/lsernau/grh-lean/blob/72fead102f663d5c7a407964a13dc0cf9d977207/grh-compact/src/Lib/LocalGL/Green/LabelIdempotents.lean\#L39-L57}{{\leanfont labelChar\_mul\_det}}}
\expandafter\newcommand\csname LN@isCuspidalgreenTau\endcsname{\href{https://github.com/lsernau/grh-lean/blob/72fead102f663d5c7a407964a13dc0cf9d977207/grh-compact/src/Lib/LocalGL/Green/GreenTau.lean\#L26-L27}{{\leanfont isCuspidal\_greenTau}}}
\expandafter\newcommand\csname LN@existscharpolyeqpow\endcsname{\href{https://github.com/lsernau/grh-lean/blob/72fead102f663d5c7a407964a13dc0cf9d977207/grh-compact/src/Lib/LocalGL/Green/CuspidalSupport.lean\#L26-L45}{{\leanfont Hall.\allowbreak{}IsCuspidal.\allowbreak{}exists\_charpoly\_eq\_pow}}}
\expandafter\newcommand\csname LN@letotientmullog\endcsname{\href{https://github.com/lsernau/grh-lean/blob/72fead102f663d5c7a407964a13dc0cf9d977207/grh-compact/src/Lib/General/Analysis/OstrowskiTaussky.lean\#L12-L19}{{\leanfont le\_totient\_mul\_log}}}
\newcommand{\ConcordanceRows}{\crow{IsPrimitiveLabel}{\href{https://github.com/lsernau/grh-lean/blob/72fead102f663d5c7a407964a13dc0cf9d977207/grh-compact/src/New/Trace/BadShare.lean\#L29-L30}{{\leanfont IsPrimitiveLabel}}}{BadShare.lean}{29--30}
\crow{scaleRoots}{\href{https://github.com/lsernau/grh-lean/blob/72fead102f663d5c7a407964a13dc0cf9d977207/grh-compact/src/New/Trace/BadShare.lean\#L35-L41}{{\leanfont IsPrimitiveLabel.\allowbreak{}scaleRoots}}}{BadShare.lean}{35--41}
\crow{isPrimitiveRootprodroots}{\href{https://github.com/lsernau/grh-lean/blob/72fead102f663d5c7a407964a13dc0cf9d977207/grh-compact/src/New/Trace/BadShare.lean\#L45-L51}{{\leanfont IsPrimitiveLabel.\allowbreak{}isPrimitiveRoot\_prod\_roots}}}{BadShare.lean}{45--51}
\crow{goodTuples}{\href{https://github.com/lsernau/grh-lean/blob/72fead102f663d5c7a407964a13dc0cf9d977207/grh-compact/src/New/Trace/BadShare.lean\#L55-L57}{{\leanfont goodTuples}}}{BadShare.lean}{55--57}
\crow{existslabelIdemeqoffirreducible}{\href{https://github.com/lsernau/grh-lean/blob/72fead102f663d5c7a407964a13dc0cf9d977207/grh-compact/src/New/Trace/BadShare.lean\#L64-L86}{{\leanfont exists\_labelIdem\_eq\_off\_irreducible}}}{BadShare.lean}{64--86}
\crow{badsharele}{\href{https://github.com/lsernau/grh-lean/blob/72fead102f663d5c7a407964a13dc0cf9d977207/grh-compact/src/New/Trace/BadShare.lean\#L92-L136}{{\leanfont bad\_share\_le}}}{BadShare.lean}{92--136}
\crow{IsLabel}{\href{https://github.com/lsernau/grh-lean/blob/72fead102f663d5c7a407964a13dc0cf9d977207/grh-compact/src/Lib/LocalGL/Green/FrobeniusOrbits.lean\#L48-L48}{{\leanfont IsLabel}}}{FrobeniusOrbits.lean}{48--48}
\crow{labelIdem}{\href{https://github.com/lsernau/grh-lean/blob/72fead102f663d5c7a407964a13dc0cf9d977207/grh-compact/src/Lib/LocalGL/Green/FrobeniusOrbits.lean\#L61-L63}{{\leanfont labelIdem}}}{FrobeniusOrbits.lean}{61--63}
\crow{tupleIdem}{\href{https://github.com/lsernau/grh-lean/blob/72fead102f663d5c7a407964a13dc0cf9d977207/grh-compact/src/Lib/LocalGL/Green/LabelIdempotents.lean\#L108-L108}{{\leanfont tupleIdem}}}{LabelIdempotents.lean}{108--108}
\crow{isCentralProjsumtupleIdem}{\href{https://github.com/lsernau/grh-lean/blob/72fead102f663d5c7a407964a13dc0cf9d977207/grh-compact/src/Lib/LocalGL/Green/LabelIdempotents.lean\#L124-L125}{{\leanfont isCentralProj\_sum\_tupleIdem}}}{LabelIdempotents.lean}{124--125}
\crow{labelChardegreebounds}{\href{https://github.com/lsernau/grh-lean/blob/72fead102f663d5c7a407964a13dc0cf9d977207/grh-compact/src/Lib/LocalGL/Green/LabelDegree.lean\#L35-L57}{{\leanfont labelChar\_degree\_bounds}}}{LabelDegree.lean}{35--57}
\crow{labelCharmuldet}{\href{https://github.com/lsernau/grh-lean/blob/72fead102f663d5c7a407964a13dc0cf9d977207/grh-compact/src/Lib/LocalGL/Green/LabelIdempotents.lean\#L39-L57}{{\leanfont labelChar\_mul\_det}}}{LabelIdempotents.lean}{39--57}
\crow{isCuspidalgreenTau}{\href{https://github.com/lsernau/grh-lean/blob/72fead102f663d5c7a407964a13dc0cf9d977207/grh-compact/src/Lib/LocalGL/Green/GreenTau.lean\#L26-L27}{{\leanfont isCuspidal\_greenTau}}}{GreenTau.lean}{26--27}
\crow{existscharpolyeqpow}{\href{https://github.com/lsernau/grh-lean/blob/72fead102f663d5c7a407964a13dc0cf9d977207/grh-compact/src/Lib/LocalGL/Green/CuspidalSupport.lean\#L26-L45}{{\leanfont Hall.\allowbreak{}IsCuspidal.\allowbreak{}exists\_charpoly\_eq\_pow}}}{CuspidalSupport.lean}{26--45}
\crow{letotientmullog}{\href{https://github.com/lsernau/grh-lean/blob/72fead102f663d5c7a407964a13dc0cf9d977207/grh-compact/src/Lib/General/Analysis/OstrowskiTaussky.lean\#L12-L19}{{\leanfont le\_totient\_mul\_log}}}{OstrowskiTaussky.lean}{12--19}}
% the notations of the defined functions: one slot for each explicit argument of the Lean declaration
\newcommand{\IsPrimitiveLabel}[2]{$#2$ is a primitive label of degree~$#1$}
\newcommand{\goodTuples}[2]{\mathrm{Good}_{#1,#2}}
\newcommand{\IsLabel}[2]{$#2$ is a label of degree~$#1$}
\newcommand{\labelIdem}[2]{e_{#1,#2}}
\newcommand{\tupleIdem}[2]{e^{\otimes}_{#1,#2}}
\newcommand{\UnitMembers}{46}
\newcommand{\UnitInNote}{6}
\newcommand{\UnitProofs}{16}
\newcommand{\UnitSteps}{99}
\newcommand{\UnitAllSteps}{126}


\linespread{1.08}                       % 11 pt on 14.7 pt
\setlength{\parindent}{0pt}
\setlength{\parskip}{5.5pt plus 1.5pt minus 0.5pt}
\setlength{\abovedisplayskip}{7.5pt plus 2pt minus 3pt}
\setlength{\belowdisplayskip}{7.5pt plus 2pt minus 3pt}
\setlength{\abovedisplayshortskip}{3pt plus 2pt}
\setlength{\belowdisplayshortskip}{6pt plus 2pt minus 2pt}
\emergencystretch=1.5em
\clubpenalty=10000 \widowpenalty=10000 \displaywidowpenalty=10000
\relpenalty=7000 \binoppenalty=9000
\raggedbottom
\pagestyle{fancy}\fancyhf{}\renewcommand{\headrulewidth}{0pt}
\fancyfoot[C]{\sffamily\fontsize{7.6}{9}\selectfont\color{inkfaint}\thepage\,/\,\pageref*{LastPage}}
\setlist{topsep=2pt,partopsep=0pt,parsep=0pt,itemsep=4pt,leftmargin=2.1em,labelsep=0.55em}

\newcommand{\Q}{\mathbb{Q}}
\newcommand{\Z}{\mathbb{Z}}
\newcommand{\F}{\mathbb{F}}
\newcommand{\len}{\operatorname{len}}
\newcommand{\cok}{\operatorname{cok}}
\newcommand{\End}{\operatorname{End}}
\newcommand{\Hom}{\operatorname{Hom}}
\newcommand{\Gal}{\operatorname{Gal}}
\newcommand{\diag}{\operatorname{diag}}
\newcommand{\kw}[1]{\emph{#1}}
\newcommand{\m}[1]{\mbox{$#1$}}        % a short formula that must not be broken across lines
\newcommand{\head}[1]{\par\addvspace{17pt plus 4pt minus 2pt}\Needspace{6\baselineskip}%
  \noindent{\fontsize{13.2}{16}\selectfont\bfseries\color{ink}#1}\par\nopagebreak\vspace{-1pt}}
\newcommand{\sect}[2]{\refstepcounter{section}\label{#2}\par\addvspace{19pt plus 4pt minus 2pt}\Needspace{7\baselineskip}%
  \noindent\llap{{\fontsize{13.2}{16}\selectfont\bfseries\color{inkfaint}\thesection}\hspace{9pt}}%
  {\fontsize{13.2}{16}\selectfont\bfseries\color{ink}#1}\par\nopagebreak\vspace{-1pt}}
\newcommand{\subhead}[1]{\par\addvspace{9pt plus 2pt}\Needspace{4\baselineskip}\noindent\textbf{#1}\enspace\ignorespaces}
% where a statement lives in the Lean archive
\newcommand{\leantag}{{\sffamily\bfseries\fontsize{6.3}{8}\selectfont\color{leancol}\addfontfeatures{LetterSpace=7}LEAN}}
\newcommand{\sep}{\unskip\nobreak\hspace{5pt}{\color{inkfaint}\textperiodcentered}\hspace{5pt plus 2pt}\ignorespaces}
\newcommand{\leanline}[1]{\par\nopagebreak\vspace{-2.5pt}{\fontsize{8.9}{12.4}\selectfont\color{inksoft}\raggedright
  \noindent\llap{\leantag\hspace{9pt}}#1\par}\vspace{1pt}}
% a displayed statement in a tinted box: #1 colour, #2 text
\newcommand{\statement}[2]{\par\addvspace{7pt}\begingroup\setlength{\fboxsep}{8.5pt}\setlength{\fboxrule}{0.6pt}%
  \noindent\fcolorbox{#1!30!white}{#1!3!white}{\parbox{\dimexpr\linewidth-2\fboxsep-2\fboxrule\relax}{#2}}%
  \endgroup\par\addvspace{7pt}}
\newcommand{\nt}[2]{\par\noindent\hangindent=1.1em\hangafter=1 #1\enspace{\color{inksoft}#2}\par\vspace{1pt}}
\newcommand{\legdot}[1]{\tikz[baseline=-0.62ex]\fill[#1] (0,0) circle (2.5pt);\hspace{3pt}}
\newcommand{\crow}[4]{\csname CS@#1\endcsname & #2 & \textcolor{inksoft}{#3} & \textcolor{inksoft}{#4}\\}
\newcommand{\leanquoteline}[2]{\par\noindent\hangindent=\dimexpr#1\fontdimen2\font+2.2em\relax\hangafter=1
  \hspace*{#1\fontdimen2\font}#2}
\newcommand{\quoted}[1]{\par\vspace{3pt}{\ttfamily\fontsize{7.5}{10.4}\selectfont\color{leancol}\raggedright
  \setlength{\parskip}{0pt}\leanquote{#1}\par}\vspace{4pt}}
\newcommand{\bx}[3]{{\sffamily\bfseries\fontsize{7.3}{9.2}\selectfont\color{#1!82!black}#2}\\[1.5pt]#3}
\newcommand{\cs}[2]{\expandafter\newcommand\csname CS@#1\endcsname{#2}}

% theorems
\newtheoremstyle{plainx}{7pt plus 2pt}{7pt plus 2pt}{\itshape}{}{\bfseries}{.}{0.5em}{}
\newtheoremstyle{uprightx}{7pt plus 2pt}{7pt plus 2pt}{}{}{\bfseries}{.}{0.5em}{}
\theoremstyle{plainx}
\newtheorem{theorem}{Theorem}[section]
\newtheorem{lemma}[theorem]{Lemma}
\newtheorem{proposition}[theorem]{Proposition}
\newtheorem{corollary}[theorem]{Corollary}
\theoremstyle{uprightx}
\newtheorem{definition}[theorem]{Definition}
\newtheorem{setting}[theorem]{Setting}          % hypotheses shared by several statements; each of them cites it by number
\newtheorem{remark}[theorem]{Remark}
\newtheorem{example}[theorem]{Example}
\renewcommand{\qedsymbol}{{\setlength{\fboxsep}{0pt}\setlength{\fboxrule}{0.5pt}\fbox{\rule{0pt}{0.46em}\hspace{0.46em}}}}
% what holds throughout the note: the ambient objects only (implicit and instance arguments), never a hypothesis
\newenvironment{conventions}{\par\addvspace{5pt}\noindent\textbf{Conventions.}\enspace\ignorespaces}{\par\addvspace{3pt}}
\makeatletter
\renewenvironment{proof}[1][Proof]{\par\pushQED{\qed}\normalfont\topsep0pt\partopsep0pt
  \trivlist\item[\hskip\labelsep\itshape #1.]\ignorespaces}{\popQED\endtrivlist\@endpefalse}
\makeatother

% where each declaration of the concordance is treated: \cs{key}{place}, the place naming a \label or a lettered item
\cs{IsPrimitiveLabel}{Definition~\ref{d:prim}}
\cs{scaleRoots}{Lemma~\ref{l:scale}}
\cs{isPrimitiveRootprodroots}{Lemma~\ref{l:norm}}
\cs{goodTuples}{Definition~\ref{d:good}}
\cs{existslabelIdemeqoffirreducible}{Theorem 2}
\cs{badsharele}{Theorem 1}
\cs{IsLabel}{Definition~\ref{d:label}}
\cs{labelIdem}{Definition~\ref{d:idem}}
\cs{tupleIdem}{Definition~\ref{d:tuple}}
\cs{isCentralProjsumtupleIdem}{Input 1}
\cs{labelChardegreebounds}{Input 1}
\cs{labelCharmuldet}{Input 2}
\cs{isCuspidalgreenTau}{Input 2}
\cs{existscharpolyeqpow}{Input 2}
\cs{letotientmullog}{Input 3}

\begin{document}
\color{ink}

{\raggedright
{\sffamily\bfseries\fontsize{7.2}{9}\selectfont\color{leancol}\addfontfeatures{LetterSpace=8}READBACK NOTE 19\par}
\vspace{5pt}
{\fontsize{21.5}{25.5}\selectfont\bfseries Primitive labels of general linear groups over finite fields: the share of bad types and the cubic twist\par}
\vspace{5pt}
{\fontsize{12.4}{16}\selectfont\itshape\color{inksoft} One unit of the Lean artifact \ArtifactName, read back into
mathematics\par}}
\vspace{8pt}
{\color{rulegrey}\hrule height 0.6pt}
\vspace{7pt}
\thispagestyle{empty}

\textbf{Summary.} The rows of the proof are cut out of a space of automorphic forms by an idempotent of the group
algebra of a product of finite general linear groups. This note treats the part of that construction which lives
over finite fields. A \emph{label} is a squarefree monic polynomial; it names an irreducible character of
$\mathrm{GL}_r$ of a finite field, by Green's construction. A label is \emph{primitive} when it is irreducible
with a root that generates the multiplicative group of the extension of degree~$r$. Two results are proved.
Theorem~1: in the regular representation of a product of $f$ groups $\mathrm{GL}_r(\F_v)$, the characters named
by tuples of labels that avoid finitely many roots and are primitive at one factor at least have all but a share
\m{f\varepsilon+\exp(-L)}. Theorem~2: when $r$ is a power of~$3$, a primitive label has a cubic twist that avoids two
given polynomials and whose character agrees with the original one at every element with a reducible
characteristic polynomial.

\statement{leancol}{\fontsize{9.6}{13.2}\selectfont\textbf{What this note is.} The artifact \ArtifactName{} is a
body of Lean code: the compact edition of a formal proof, due to L.~Alpöge with the AI system Claude, of the
Riemann hypothesis and of its generalisation to Dirichlet $L$-functions. That headline is a claim. The authors
report that every file compiles and that Lean's kernel accepted the proofs; an independent replay is recorded in the
handout's record of checks when it reports. No mathematician has read the argument through, and nothing here says that the Riemann
hypothesis is proved. This note treats one file of the unit \emph{Good idempotents, labels and the transfer of
traces}: \UnitInNote{} of the unit's \UnitMembers{} declarations (the other declarations are the subject of the
note \emph{Good forms and rows}). The two theorems and the two lemmas are proved in full. Taken for granted,
and stated as Inputs~1 to~3, are facts about Green's characters and about Euler's function that the artifact
proves in its background library. Every statement was compared with its Lean text, argument by argument. The Lean
text read is \PinText; Appendix~A is the concordance of \ConcDecls{} declarations.}

\sect{The statements}{s:intro}
The notions are defined in §\ref{s:setting}: labels and primitive labels in Definitions~\ref{d:label}
and~\ref{d:prim}, the idempotents \m{\labelIdem{I}{c}} and \m{\tupleIdem{I}{\mathbf c}} in
Definitions~\ref{d:idem} and~\ref{d:tuple}, the set \m{\goodTuples{M}{r}} of good tuples in
Definition~\ref{d:good}. For an element $y$ of the complex group algebra of a finite group and an element $g$ of
the group, $y(g)$ is the coefficient of $y$ at~$g$.

\statement{arith}{\phantomsection\label{t:1}\textbf{Theorem 1} (bad types have a small share).\enspace
\emph{Let $I$ and $A$ be finite sets, let \m{(\F_v)_{v\in A}} be a family of finite fields, let $q$, $M$, $r$, $f$ be
natural numbers with \m{M\ge1}, and let $\varepsilon$ and $L$ be real numbers. Suppose that every $\F_v$ has $q$
elements, that $I$ has $r$ elements and $A$ has $f$ elements, and that}
\[
1<r,\qquad r<q,\qquad \varepsilon\le1,\qquad \frac{1+M+2r}{q}\le\varepsilon,\qquad
2r\,\bigl(2r\log q+1\bigr)\,L\le f .
\]
\emph{Then the coefficient at the identity of the element \m{1-\sum_{\mathbf c\in\goodTuples{M}{r}}\tupleIdem{I}{\mathbf c}}
of the complex group algebra of \m{\prod_{v\in A}\mathrm{GL}_I(\F_v)} has real part at most}
\[
f\,\varepsilon+\exp(-L).
\]}
\leanline{\lean{badsharele}}

\statement{arith}{\phantomsection\label{t:2}\textbf{Theorem 2} (the cubic twist).\enspace
\emph{Let $\F$ be a finite field, $I$ a finite set and $a$ a natural number. Suppose that $I$ has \m{3^{a+1}}
elements and that $9$ divides \m{\#\F-1}. Let \m{c\in\F[X]} be a polynomial, and let $\mathcal S_1$ and
$\mathcal S_2$ be sets of polynomials over~$\F$ with at most one element each. Then there is a polynomial
\m{\tilde c\in\F[X]} with the following two properties.}
\begin{enumerate}[label=\textup{(\arabic*)},leftmargin=2.3em,itemsep=2pt]
\item \emph{If \IsPrimitiveLabel{\#I}{c}, then \IsPrimitiveLabel{\#I}{\tilde c}, and
\m{\tilde c\notin\mathcal S_1\cup\mathcal S_2}.}
\item \emph{For every \m{g\in\mathrm{GL}_I(\F)} whose characteristic polynomial is not irreducible,
\m{\labelIdem{I}{c}(g)=\labelIdem{I}{\tilde c}(g)}.}
\end{enumerate}}
\leanline{\lean{existslabelIdemeqoffirreducible}}

In Theorem~1 the element \m{\sum_{\mathbf c}\tupleIdem{I}{\mathbf c}} is a central idempotent (Input~1), and the
coefficient at the identity of a central idempotent is the share of its image in the regular representation: the
dimension of the image divided by the order of the group. So Theorem~1 says that the representations whose type
is not a good tuple fill at most the share \m{f\varepsilon+\exp(-L)} of the group algebra. In the proof of the
main theorems the fields are the residue fields at the chosen primes above one of two auxiliary primes, $M$ is a
fixed integer, and $f$ is the number of infinite places of a large CM field; the theorem is used in
Chapter~VI. Theorem~2 is the device that makes a comparison of trace formulas possible in the note \emph{Good
forms and rows}: subtracting from an idempotent its cubic twist leaves an element of the group algebra supported
on matrices with irreducible characteristic polynomial.

\sect{The setting}{s:setting}
\begin{conventions}\phantomsection\label{c:conventions}
The natural numbers are \m{0,1,2,\dots}. Throughout, $\F$ is a finite field, $\bar{\F}$ an algebraic closure
of~$\F$, and \m{\#\F} the number of elements of~$\F$; in a family \m{(\F_v)_{v\in A}} of finite fields the index
set $A$ is a finite set. A finite set $I$ serves as index set of matrices: \m{\mathrm{GL}_I(\F)} is the group of
invertible matrices over~$\F$ with rows and columns indexed by~$I$, and \m{\#I} is the number of elements of~$I$.
The letters $r$, $M$, $q$, $f$, $a$ denote natural numbers and $c$ a polynomial in \m{\F[X]}. For a finite group
$G$, \m{\mathbb C[G]} is its complex group algebra, $y(g)$ is the coefficient of \m{y\in\mathbb C[G]} at \m{g\in G},
and a complex representation of~$G$ is also a module over \m{\mathbb C[G]}. A multiplicative character of~$\F$ is
a homomorphism \m{\chi\colon\F^\times\to\mathbb C^\times}. For a unit \m{u\in\F^\times} and a polynomial \m{c\in\F[X]}
of degree~$\delta$, \m{c^{[u]}=u^{\delta}\,c(X/u)} is the polynomial of the same degree whose roots in $\bar{\F}$ are
$u$ times the roots of~$c$, with the same multiplicities; \m{c\mapsto c^{[u]}} is a bijection of \m{\F[X]} that
respects products, with inverse \m{c\mapsto c^{[u^{-1}]}}. Letters of the handout's overview that are not
displayed with the data as an argument (among them $\varepsilon$, $\delta$, $f$, $q$, $n$ and~$V$) have, in this
note, only the meaning given in its notation list.
\end{conventions}

\begin{definition}[labels]\label{d:label}
Let $r$ be a natural number and \m{c\in\F[X]}. Then \IsLabel{r}{c} if $c$ is monic and squarefree, has degree~$r$,
and has \m{c(0)\ne0}.
\end{definition}
\leanline{\lean{IsLabel}\sep a definition of the background library}

\begin{definition}[primitive labels]\label{d:prim}
Let $r$ be a natural number and \m{c\in\F[X]}. Then \IsPrimitiveLabel{r}{c} if \IsLabel{r}{c}, if $c$ is
irreducible, and if some root \m{\xi\in\bar{\F}} of~$c$ is a primitive root of unity of order \m{(\#\F)^r-1}.
\end{definition}
\leanline{\lean{IsPrimitiveLabel}}

So a primitive label of degree~$r$ is the minimal polynomial over~$\F$ of a generator of the cyclic group
\m{\F_{(r)}^\times}, where \m{\F_{(r)}\subset\bar{\F}} is the subfield with \m{(\#\F)^r} elements. All roots of a
primitive label are such generators, because they are conjugate over~$\F$.

\begin{definition}[Green's characters and their idempotents]\label{d:idem}
Fix an injective multiplicative character \m{T\colon\bar{\F}^\times\to\mathbb C^\times} (the library chooses one for
each field). For an integer \m{\delta\ge1} and \m{x\in\F_{(\delta)}^\times}, let $\theta_x$ be the character
$T^{n}$ of \m{\F_{(\delta)}^\times}, where $n$ is an integer with \m{T(x)=\exp\bigl(2\pi i\,n/((\#\F)^\delta-1)\bigr)};
the map \m{x\mapsto\theta_x} is an isomorphism from \m{\F_{(\delta)}^\times} to its group of characters. For
\m{\delta=1} it gives, for each multiplicative character $\chi$ of~$\F$, the unit \m{x_\chi\in\F^\times} with
\m{\theta_{x_\chi}=\chi}, and \m{\chi\mapsto x_\chi} is an isomorphism of groups.

The complex functions on the similarity classes of square matrices over~$\F$ of all sizes form a ring under the
product of parabolic induction,
\[
(F*H)(g)=\sum_W F(g|_W)\,H(g|_{V/W}),
\]
the sum over the $g$-stable subspaces $W$ of the space $V$ on which $g$ acts. For a monic irreducible \m{P\in\F[X]} of degree~$\delta$ with a root
\m{x\in\bar{\F}^\times}, Green's class function $\tau_P$ is \m{(-1)^\delta} times the part in size~$\delta$ of
\m{(\beta_x\cdot\mathbb 1^{-1})*\mathbb 1}; here $\mathbb 1$ is the constant function~$1$, \m{\mathbb 1^{-1}} is its inverse in this
ring, the dot is the pointwise product, and
\m{\beta_x(g)} is the sum of $\theta_x$ over the eigenvalues of~$g$. For a polynomial \m{c\ne0} over~$\F$,
\[
\Xi_c=\tau_{P_1}*\dots*\tau_{P_s},\qquad P_1,\dots,P_s\ \text{the monic irreducible factors of $c$ with multiplicity.}
\]
For a finite set~$I$ and \m{c\in\F[X]}, the element \m{\labelIdem{I}{c}} of \m{\mathbb C[\mathrm{GL}_I(\F)]} is
\[
\labelIdem{I}{c}=\sum_{g\in\mathrm{GL}_I(\F)}\frac{\Xi_c(1)}{\#\mathrm{GL}_I(\F)}\;\overline{\Xi_c(g)}\;g
\]
if \IsLabel{\#I}{c}, and \m{\labelIdem{I}{c}=0} otherwise.
\end{definition}
\leanline{\lean{labelIdem}\sep a definition of the background library, with the class functions of its
files on Green's characters}

For a label~$c$ of degree \m{\#I} the restriction of $\Xi_c$ to \m{\mathrm{GL}_I(\F)} is an irreducible
character, by Green's theory as the library proves it; then \m{\labelIdem{I}{c}} is the central idempotent of
\m{\mathbb C[\mathrm{GL}_I(\F)]} that projects every representation to its part of type~$\Xi_c$. If $c$ is
irreducible, \m{\Xi_c=\tau_c} is a cuspidal character. This note uses nothing of the construction beyond
Inputs~1 to~3.

\begin{definition}[idempotents of tuples]\label{d:tuple}
Let \m{(\F_v)_{v\in A}} be a family of finite fields, $I$ a finite set, and \m{\mathbf c=(c_v)_{v\in A}} a family of
polynomials \m{c_v\in\F_v[X]}. Then \m{\tupleIdem{I}{\mathbf c}} is the element
\m{\bigotimes_{v\in A}\labelIdem{I}{c_v}} of the group algebra of \m{\prod_{v\in A}\mathrm{GL}_I(\F_v)}: its
coefficient at \m{(g_v)_{v\in A}} is \m{\prod_{v\in A}\labelIdem{I}{c_v}(g_v)}.
\end{definition}
\leanline{\lean{tupleIdem}\sep a definition of the background library}

\begin{definition}[good tuples]\label{d:good}
Let \m{(\F_v)_{v\in A}} be a family of finite fields, and let $M$ and $r$ be natural numbers. The set
\m{\goodTuples{M}{r}} consists of the families \m{\mathbf c=(c_v)_{v\in A}} of monic polynomials \m{c_v\in\F_v[X]}
of degree~$r$ such that
\begin{enumerate}[label=\textup{(\roman*)},leftmargin=2.3em,itemsep=2pt]
\item for every \m{v\in A}: \IsLabel{r}{c_v}, and no \m{x\in\F_v} with \m{x^M=1} is a root of~$c_v$;
\item for at least one \m{v\in A}: \IsPrimitiveLabel{r}{c_v}.
\end{enumerate}
It is a finite set.
\end{definition}
\leanline{\lean{goodTuples}}

In terms of characters, (i) says that the type \m{\Xi_{c_v}} is induced from cuspidal characters among which
there is no character \m{\chi\circ\det} of \m{\mathrm{GL}_1} with \m{\chi^M=1}, and (ii) asks for one factor
with a cuspidal type that comes from a faithful character of \m{\F_{v,(r)}^\times}.

\statement{inksoft}{\textbf{Inputs from the background library.} The statements below are proved in the
artifact's library and are taken for granted here. For Inputs 1 and~2 the library's files name
J.~A.~Green, \emph{The characters of the finite general linear groups}, Trans. Amer. Math. Soc. 80 (1955),
as their source. The second reading of this write-up compared Input~1 with that paper: both parts are weaker
consequences of its §8, Theorems 13 and~14 (pp.~441--444, with the degree formula on pp.~443--444). For Input~2
the place in the paper is to be confirmed; the catalogue of quoted results classes Input~2(b) as folklore.

\vspace{4pt}\phantomsection\label{i:1}\textbf{Input 1} (orthogonality and degree).\enspace
\emph{Let $I$ be a finite set and \m{(\F_v)_{v\in A}} a family of finite fields. \textup{(a)} For every finite set
$\mathcal C$ of families \m{\mathbf c=(c_v)_{v\in A}} of polynomials \m{c_v\in\F_v[X]}, the element
\m{\sum_{\mathbf c\in\mathcal C}\tupleIdem{I}{\mathbf c}} is an idempotent whose coefficients are constant on
conjugacy classes. \textup{(b)} If $\F$ has $q$ elements, $I$ has $r$ elements and \IsLabel{r}{c}, then the value
\m{\Xi_c(1)} at the identity of \m{\mathrm{GL}_I(\F)} has real part \m{\operatorname{Re}\Xi_c(1)\ge0}, and
\m{\labelIdem{I}{c}(1)=\Xi_c(1)\,\overline{\Xi_c(1)}/\#\mathrm{GL}_I(\F)} is a real number with
\m{\labelIdem{I}{c}(1)\ge(1-r/q)\,q^{-r}}.}

\vspace{3pt}\phantomsection\label{i:2}\textbf{Input 2} (twisting; cuspidal support).\enspace
\emph{\textup{(a)} Let $r$ be a natural number and suppose that \IsLabel{r}{c}; let $\chi$ be a multiplicative
character of~$\F$ and $I$ a finite set. Then \m{\Xi_c(g)\,\chi(\det g)=\Xi_{c^{[x_\chi]}}(g)} for all
\m{g\in\mathrm{GL}_I(\F)}. \textup{(b)} Let \m{P\in\F[X]} be monic and irreducible, and let $g$ be a square matrix
over~$\F$ with \m{\tau_P(g)\ne0}. Then the characteristic polynomial of~$g$ is a power of an irreducible
polynomial.}

\vspace{3pt}\phantomsection\label{i:3}\textbf{Input 3} (Euler's function).\enspace
\emph{For every integer \m{n\ge1}: \m{n\le\varphi(n)\,(2\log n+1)}, where $\varphi$ is Euler's function.}}
\leanline{\lean{isCentralProjsumtupleIdem}\sep \lean{labelChardegreebounds}\sep \lean{labelCharmuldet}\sep
\lean{isCuspidalgreenTau}\sep \lean{existscharpolyeqpow}\sep \lean{letotientmullog}}

Input~1(b) is a weak form of Green's formula for the degree of~$\Xi_c$: only this lower bound is proved, and only
for labels. (As the second reading records it: for a label whose irreducible factors have degrees \m{d_1,\dots,d_s},
Green's degree is \m{\prod_{n=1}^{r}(q^n-1)/\prod_{i=1}^{s}(q^{d_i}-1)} up to sign, and the bound follows
from \m{\prod_i(q^{d_i}-1)\le q^r}. The library does prove that the real part of \m{\Xi_c(1)} is not negative; the note
uses only the bound for \m{\labelIdem{I}{c}(1)}.) Input~2(b) is the vanishing of cuspidal characters off the primary classes; the library proves it
for every class function that is cuspidal in the sense that its averages over block triangular matrices with two
diagonal blocks of positive size vanish, and proves that $\tau_P$ is cuspidal in this sense.

\sect{Primitive labels}{s:prim}
Two facts about finite fields are used without further mention. If \m{P\in\F[X]} is irreducible of degree~$\delta$
and \m{x\in\bar{\F}} is a root of~$P$, then the roots of~$P$ are \m{x,x^{q},\dots,x^{q^{\delta-1}}} with \m{q=\#\F},
each a simple root; and \m{x^{q^n}=x} holds exactly when $\delta$ divides~$n$.

\begin{lemma}[scaling the roots]\label{l:scale}
Let \m{c\in\F[X]}, let $u$ be a unit of~$\F$, and let $r$ and $M$ be natural numbers. Suppose that \m{u^M=1} and
that $M^2$ divides \m{\#\F-1}. If \IsPrimitiveLabel{r}{c}, then \IsPrimitiveLabel{r}{c^{[u]}}.
\end{lemma}
\begin{proof}
Put \m{q=\#\F}. The polynomial \m{c^{[u]}} is monic of degree~$r$, its roots are $u$ times those of~$c$, hence
distinct, and \m{c^{[u]}(0)=u^r c(0)\ne0}; so \IsLabel{r}{c^{[u]}}. It is irreducible, because
\m{c\mapsto c^{[u]}} is a bijection of \m{\F[X]} that respects products and degrees. Let $\xi$ be a root of~$c$ of
order \m{n=q^r-1}; then $u\xi$ is a root of \m{c^{[u]}}, and it remains to see that $u\xi$ has order~$n$. Let $o$
be the order of~$u$. It divides~$M$, so $o^2$ divides \m{q-1}, which divides~$n$. Since \m{(u\xi)^n=1}, it
suffices to show \m{(u\xi)^{n/p}\ne1} for every prime~$p$ dividing~$n$. If $p$ divides~$o$, then $po$ divides
$o^2$ and so~$n$; hence \m{u^{n/p}=1} and \m{(u\xi)^{n/p}=\xi^{n/p}\ne1}. If $p$ does not divide~$o$ and
\m{(u\xi)^{n/p}=1}, then \m{\xi^{n/p}=u^{-n/p}}; the left side has order~$p$ and the order of the right side
divides~$o$, a contradiction.
\end{proof}
\leanline{\lean{scaleRoots}}

\begin{lemma}[the norm of a generator is a generator]\label{l:norm}
Let \m{c\in\F[X]} and let $r$ be a natural number. Suppose \IsPrimitiveLabel{r}{c}. Then the product of the
roots of~$c$ in~$\bar{\F}$, counted with multiplicity, is a primitive root of unity of order \m{\#\F-1}; and there
is a natural number~$n$ such that \m{\xi^{n}} equals this product for every root $\xi$ of~$c$ in~$\bar{\F}$.
\end{lemma}
\begin{proof}
Put \m{q=\#\F} and \m{n=1+q+\dots+q^{r-1}}. Let $\xi$ be any root of~$c$. As $c$ is irreducible of degree~$r$, its
roots are \m{\xi^{q^i}}, \m{0\le i<r}, each once, and their product is~$\xi^{n}$. For the root $\xi_0$ of order
\m{q^r-1=n\,(q-1)} the power \m{\xi_0^{n}} has order \m{q-1}.
\end{proof}
\leanline{\lean{isPrimitiveRootprodroots}}

\sect{The cubic twist: proof of Theorem 2}{s:twist}
\begin{proof}[Proof of Theorem 2]
Put \m{q=\#\F} and \m{r=\#I=3^{a+1}}. If $c$ is not a primitive label of degree~$r$, take \m{\tilde c=c}: then
(1) is empty and (2) is trivial. So assume that \IsPrimitiveLabel{r}{c}, and let $\xi$ be a root of~$c$ of order \m{q^r-1}.

\emph{The element $\xi^3$ has degree $r$ over $\F$.} Its degree $\delta$ divides~$r$, since
\m{\xi^3\in\F_{(r)}}. If \m{\delta\ne r}, then $\delta$ divides \m{r/3}, so $\xi^3$ lies in the field with
\m{z=q^{r/3}} elements and its order \m{(q^r-1)/3} divides \m{z-1}. Then \m{z^3-1\le3(z-1)}, that is
\m{(z-1)^2(z+2)\le0}, which is false for \m{z\ge2}. (The order of $\xi^3$ is \m{(q^r-1)/3} because $3$ divides
\m{q-1}.)

\emph{The three cubic twists of $c$ are distinct.} As $3$ divides \m{q-1}, the group $U_3$ of the \m{u\in\F} with
\m{u^3=1} has three elements. Let \m{u,u'\in U_3} with \m{c^{[u]}=c^{[u']}}. This polynomial is irreducible and has
the roots $u\xi$ and $u'\xi$, so \m{u'\xi=(u\xi)^{q^l}=u\,\xi^{q^l}} for some $l$ with \m{0\le l<r}. Taking cubes,
\m{\xi^3=(\xi^3)^{q^l}}; since $\xi^3$ has degree~$r$, $r$ divides~$l$, so \m{l=0} and \m{u'=u}.

\emph{Choice of $\tilde c$.} Of the three distinct polynomials \m{c^{[u]}}, \m{u\in U_3}, at most two lie in
\m{\mathcal S_1\cup\mathcal S_2}. Choose \m{u\in U_3} with \m{c^{[u]}\notin\mathcal S_1\cup\mathcal S_2}, and put
\m{\tilde c=c^{[u]}}. By Lemma~\ref{l:scale} with \m{M=3}, which applies because \m{u^3=1} and $9$ divides
\m{q-1}, \IsPrimitiveLabel{r}{\tilde c}. This is (1).

\emph{The coefficients.} Let $\chi$ be the multiplicative character of~$\F$ with \m{x_\chi=u}; then \m{\chi^3=1},
because \m{\chi\mapsto x_\chi} is an isomorphism of groups. By Input~2(a),
\m{\Xi_{\tilde c}(g)=\Xi_c(g)\,\chi(\det g)} for \m{g\in\mathrm{GL}_I(\F)}, in particular
\m{\Xi_{\tilde c}(1)=\Xi_c(1)}. Both $c$ and $\tilde c$ are labels of degree~$r$, so by
Definition~\ref{d:idem}
\[
\labelIdem{I}{\tilde c}(g)=\labelIdem{I}{c}(g)\cdot\overline{\chi(\det g)}\qquad(g\in\mathrm{GL}_I(\F)).
\]
Let $g$ have a characteristic polynomial that is not irreducible. If \m{\Xi_c(g)=0}, both coefficients vanish.
Otherwise, $c$ being irreducible, \m{\Xi_c=\tau_c}, and by Input~2(b) the characteristic polynomial of~$g$ is
$P^{n}$ with $P$ irreducible and \m{n\ge1}. Comparing degrees, $n$ divides \m{3^{a+1}}; and \m{n\ne1}, because
the characteristic polynomial is not irreducible. So $3$ divides~$n$, and, $r$ being odd,
\[
\det g=(-1)^r\,P(0)^{n}=\bigl(-P(0)^{n/3}\bigr)^3
\]
is a cube in~$\F^\times$. Hence \m{\chi(\det g)=1}, and the two coefficients agree. This is (2).
\end{proof}

\sect{The share of bad types: proof of Theorem 1}{s:share}
\begin{proof}[Proof of Theorem 1]
If \m{f=0}, then $A$ is empty, the only family is the empty one, and it does not satisfy (ii) of
Definition~\ref{d:good}; so the element in question is~$1$, with coefficient~$1$ at the identity, and
\m{L\le0} by the last hypothesis, so \m{\exp(-L)\ge1}. Let \m{f\ge1}. Put
\[
\varepsilon_0=\Bigl(1-\frac rq\Bigr)q^{-r},\qquad\text{so that}\quad 0\le\varepsilon_0\le q^{-r}.
\]
For \m{v\in A} let $\mathcal P_v$ be the set of the monic polynomials of degree~$r$ over~$\F_v$; it has $q^r$
elements. Let \m{\mathcal B_v\subseteq\mathcal P_v} be the set of those that are not primitive labels of
degree~$r$, and let \m{\mathcal A_v\subseteq\mathcal P_v} be the set of those that are neither of the form
\m{(X-x)\,h} with \m{x\in\F_v}, \m{x=0} or \m{x^M=1}, and $h$ monic of degree \m{r-1}, nor of the form
\m{u^2h} with $u$ monic of positive degree and $h$ monic.

\emph{Step 1: \m{\#\mathcal A_v\cdot\varepsilon_0\ge1-\varepsilon}.} There are at most \m{1+M} elements $x$ and
\m{q^{r-1}} polynomials~$h$, hence at most \m{(1+M)\,q^{r-1}} polynomials of the first form. A polynomial of
degree~$r$ of the second form is given by the degree \m{d<r} of~$h$, with \m{r-d} even, by $h$ and by~$u$; for each
$d$ there are at most \m{q^{d+(r-d)/2}\le q^{r-1}} of them, hence at most \m{r\,q^{r-1}} in all. So
\m{\#\mathcal A_v\ge q^r\bigl(1-(1+M+r)/q\bigr)}, and, as \m{\varepsilon_0\ge0},
\[
\#\mathcal A_v\cdot\varepsilon_0\ \ge\ \Bigl(1-\frac{1+M+r}{q}\Bigr)\Bigl(1-\frac rq\Bigr)\ \ge\ 1-\frac{1+M+2r}{q}\ \ge\ 1-\varepsilon .
\]

\emph{Step 2: \m{\#\mathcal B_v\cdot\varepsilon_0\le1-L/f}.} The field $\bar{\F}_v$ contains exactly
\m{\varphi(q^r-1)} primitive roots of unity of order \m{q^r-1}. Let $\xi$ be one of them and $\delta$ the degree of
its minimal polynomial~$P$ over~$\F_v$. Then \m{\xi^{q^\delta-1}=1}, so \m{q^r-1} divides \m{q^\delta-1}; and
\m{\xi^{q^r}=\xi}, so $\delta$ divides~$r$ and \m{q^\delta-1} divides \m{q^r-1}. Hence \m{\delta=r}, and $P$ is
monic, irreducible, of degree~$r$, with \m{P(0)\ne0} and with the root~$\xi$: a primitive label of degree~$r$.
Each such $P$ has $r$ roots, so there are at least \m{\varphi(q^r-1)/r} primitive labels of degree~$r$
over~$\F_v$. By Input~3, and because \m{q^r\ge2},
\[
q^r\le2\,(q^r-1)\le2\,\varphi(q^r-1)\bigl(2\log(q^r-1)+1\bigr)\le2r\,(2r\log q+1)\cdot\frac{\varphi(q^r-1)}{r}.
\]
With the hypothesis on~$L$ this gives \m{\#(\mathcal P_v\setminus\mathcal B_v)\ge q^r/(2r(2r\log q+1))\ge q^rL/f}.
Therefore \m{\#\mathcal B_v\le q^r(1-L/f)}, and \m{\#\mathcal B_v\cdot\varepsilon_0\le\#\mathcal B_v\,q^{-r}\le1-L/f}.

\emph{Step 3: a family \m{\mathbf c\in\prod_v\mathcal A_v} that does not lie in \m{\prod_v\mathcal B_v} belongs to
\m{\goodTuples{M}{r}}.} Let \m{c_v\in\mathcal A_v}. It is monic of degree~$r$. It has no root \m{x\in\F_v} with
\m{x=0} or \m{x^M=1}, since a root~$x$ gives \m{c_v=(X-x)h} with $h$ monic of degree \m{r-1}. It is squarefree:
otherwise the square of a monic irreducible polynomial~$u$ divides it, and \m{c_v=u^2h} with $h$ monic. So
\IsLabel{r}{c_v}, and (i) holds; (ii) holds because some $c_v$ is not in~$\mathcal B_v$.

\emph{Step 4: conclusion.} Let
\[
y=\sum_{\mathbf c\in\goodTuples{M}{r}}\tupleIdem{I}{\mathbf c},\qquad\text{where}\quad
\tupleIdem{I}{\mathbf c}(1)=\prod_{v\in A}\labelIdem{I}{c_v}(1)
\]
by Definition~\ref{d:tuple}. For \m{\mathbf c\in\goodTuples{M}{r}} every factor is a real number
\m{\ge\varepsilon_0\ge0} by Input~1(b). So $y(1)$ is
real, and with Step~3, then Steps 1 and~2 (all factors being non-negative, and \m{1-\varepsilon\ge0}),
\[
y(1)\ \ge\ \#\goodTuples{M}{r}\cdot\varepsilon_0^{f}\ \ge\ \prod_{v\in A}\bigl(\#\mathcal A_v\,\varepsilon_0\bigr)-\prod_{v\in A}\bigl(\#\mathcal B_v\,\varepsilon_0\bigr)
\ \ge\ (1-\varepsilon)^f-\Bigl(1-\frac Lf\Bigr)^{f}.
\]
By Bernoulli's inequality, \m{(1-\varepsilon)^f\ge1-f\varepsilon}. Moreover
\m{0\le1-L/f\le \exp(-L/f)}, so that \m{(1-L/f)^f\le \exp(-L)}. Hence
\m{1-y(1)\le f\varepsilon+\exp(-L)}.
\end{proof}

\sect{Where the results are used}{s:use}
Theorem~1 is used once, in the proof of Theorem~B of the handout (enough forms, Chapter~VI): there the fields
are the residue fields at the chosen primes above one of the primes $g_0(D)$, $g_1(D)$ of the data~$D$, with $q$ that prime, the integer $M$
is the integer $M_\omega(D)$ of the data, and \m{L=9\log b(D)}; the bound says that the good idempotent discards a
negligible share of the Poincaré series. The set \m{\goodTuples{M}{r}} enters the definition of the good
idempotent of the data (note \emph{Good forms and rows}). Theorem~2 is used in the same note, in the
comparison of the traces of Hecke operators on good forms of two unitary groups. Lemmas \ref{l:scale}
and~\ref{l:norm} are used in the note \emph{At most one primitive label for a Hecke eigensystem}: the first
to twist a primitive label by a character of order dividing $M_\omega(D)$, the second to see that the roots of a
primitive label cannot all be roots of unity of order dividing~$M_\omega(D)$.

\sect{Notes on the formal text}{s:notes}
\begin{enumerate}[label=\textup{\arabic*.},leftmargin=1.6em,itemsep=5pt]
\item \emph{Letters.} The formal text writes {\leanfont k} for the field, {\leanfont n} and {\leanfont ι} for the
index sets $I$ and~$A$, {\leanfont a} for the unit $u$ of Lemma~\ref{l:scale} and, as here, for the exponent of
Theorem~2, {\leanfont S} and {\leanfont T} for the sets $\mathcal S_1$ and~$\mathcal S_2$, {\leanfont P} for the
polynomial~$\tilde c$, and {\leanfont m} for the number~$\varepsilon_0$. Of these, \emph{k}, \emph{n} and
\emph{ι} are kept in this write-up for objects of the data, which is why the field and the index sets are
written otherwise here. The letters \emph{P}, \emph{m} and \emph{a} are letters of the data as well (the
parameters, the number of lifts, and an integer that is always written with its argument, $a(D)$); none of these
objects occurs in this note, so the irreducible polynomial~$P$ and the exponent~$a$ here are local.
\item \emph{The polynomial $c^{[u]}$.} It is the library's operation of scaling the roots, \m{\sum_i c_i\,u^{\delta-i}X^i}
for \m{c=\sum_ic_iX^i} of degree~$\delta$.
\item \emph{Theorem~2 for a polynomial that is not a primitive label.} The statement is made for every
polynomial~$c$; for such a $c$ the formal proof takes \m{\tilde c=c}, as here.
\item \emph{The distinctness of the cubic twists.} The formal proof quotes a lemma of the library: for a
prime~$l$, an exponent~$n$ and $\xi$ of order \m{q^{3^{a+1}}-1} with \m{9\mid q-1}, the conjugates of
\m{\xi^{l^n}} over~$\F$ are distinct. It is used for \m{l=3}, \m{n=1}. The two lines given here for that case are
mine; they do not use that $9$ divides \m{q-1}, only that $3$ does. The hypothesis \m{9\mid\#\F-1} of Theorem~2 is
needed for Lemma~\ref{l:scale}. The formal proof finds the twist outside \m{\mathcal S_1\cup\mathcal S_2} by
contradiction: otherwise \m{u\mapsto[c^{[u]}\in\mathcal S_1]} would embed $U_3$ into a set of two elements.
\item \emph{Lemma~\ref{l:scale}.} The formal proof quotes from Mathlib that the order of a product of two
commuting elements $u$, $\xi$ equals the order of~$\xi$ when \m{p\cdot\mathrm{ord}(u)} divides \m{\mathrm{ord}(\xi)}
for every prime $p$ dividing \m{\mathrm{ord}(u)}; the proof above includes an argument for it. The case
\m{M=0} cannot occur, since $0$ does not divide \m{\#\F-1}.
\item \emph{Theorem~1.} The formal statement asks \m{M\ne0} through an instance argument. The sets
$\mathcal A_v$ are written there as differences of finite sets: the monic polynomials of degree~$r$, less the
products \m{(X-x)h} and less the products \m{u^2h} with $h$ monic of degree \m{d<r} and $u$ monic of degree
\m{\lfloor(r-d)/2\rfloor}; products of the wrong degree are removed with them, which changes nothing. The
case \m{f=0} needs no separate treatment there, because division by zero gives zero in Lean. The inequalities of
Steps 1, 2 and~4 are established in the formal text by chains of rewriting under the order relation, closed by
automatic tactics; I derived each of them by hand, as written above, and compared the terms of the chains with it.
What the automatic tactics prove in between was not replayed.
\item \emph{What Theorem~1 does not use.} Of Green's theory it uses only the lower bound of Input~1(b) for the
coefficient at the identity. It does not use that the elements \m{\tupleIdem{I}{\mathbf c}} are orthogonal
idempotents; that is needed for the interpretation as a share.
\item \emph{Arguments.} In Definition~\ref{d:prim} the formal text does not ask that $\F$ be finite; the number
\m{\#\F} is $0$ for an infinite field there, and the notion is used for finite fields only. The dossier of the
unit records the arguments of the declarations of Definition~\ref{d:good} and of Theorem~1 as derived by hand
from the text of their signatures, not as printed by Lean; they agree with what is stated here.
\end{enumerate}

\subhead{How this note was written.} From the dossier of the unit: its statements and the structured proofs of
the two lemmas and the two theorems, which a program compared line by line with the source (\UnitSteps{} leaf
steps in the \UnitProofs{} proofs of the unit). I read the whole Lean file of these six declarations, and in the
library the definitions of labels, of Green's class functions and of the idempotents, and the statements of the
six results quoted as Inputs. The statements of this note were compared with the Lean statements argument by
argument, and the order of the quantifiers is that of the formal text. The proofs follow the formal ones, with
the differences listed in notes 4 to~6. The interpretations after the statements and after
Definition~\ref{d:good} are mine. The reference and its place are those that the library's files, the catalogue of this
handout and its second reading give; I did not check them in the literature.

\head{Appendix A. Concordance with the Lean artifact}
{\fontsize{8.6}{11.6}\selectfont\renewcommand{\leanfont}{\ttfamily\fontsize{7.7}{10}\selectfont\color{leancol}}
\setlength{\LTleft}{0pt}\setlength{\LTright}{\fill}\setlength{\tabcolsep}{0pt}
\begin{longtable}{@{}p{30mm}p{56mm}p{43mm}r@{}}
\textbf{In this note} & \textbf{Lean declaration} & \textbf{File} & \textbf{Lines}\\\midrule\endhead
\ConcordanceRows
\end{longtable}}

\subhead{The two theorems as they are written in the artifact.}
\quoted{badsharele}
\quoted{existslabelIdemeqoffirreducible}
In these texts {\leanfont k} is the field or the family of fields, {\leanfont n} and {\leanfont ι} are the index
sets $I$ and~$A$, {\leanfont labelIdem n c} is \m{\labelIdem{I}{c}}, {\leanfont tupleIdem n c} is
\m{\tupleIdem{I}{\mathbf c}}, {\leanfont goodTuples M r} is \m{\goodTuples{M}{r}}, {\leanfont .card} is the number
of elements, and {\leanfont coeff} takes the coefficient of an element of a group algebra.

\head{Notation}
\begingroup\fontsize{8.6}{10.9}\selectfont\raggedright
\setlength{\columnsep}{9mm}\setlength{\parskip}{0pt}
\begin{multicols}{2}
\nt{$\F$, $\bar{\F}$, $\#\F$}{a finite field, an algebraic closure, the number of its elements}
\nt{$\F_{(\delta)}$}{the subfield of $\bar{\F}$ with \m{(\#\F)^\delta} elements}
\nt{$q$}{the number of elements of the finite field or fields at hand}
\nt{$A$, $v$, $\F_v$, $f$}{a finite index set, an index, the field with index~$v$; \m{f=\#A}}
\nt{$I$, $\mathrm{GL}_I(\F)$}{a finite set; the invertible matrices over~$\F$ indexed by~$I$}
\nt{$r$}{a degree; in the theorems \m{r=\#I}}
\nt{$c$, $\tilde c$, $P$, $h$, $u$}{polynomials over a finite field; $u$ also a unit of~$\F$}
\nt{$X$}{the variable of polynomials}
\nt{$c^{[u]}$}{the polynomial whose roots are $u$ times those of~$c$}
\nt{$\mathbf c=(c_v)_{v\in A}$}{a family of polynomials, \m{c_v\in\F_v[X]}}
\nt{$M$}{a natural number: roots $x$ with \m{x^M=1} are excluded}
\nt{$\varepsilon$, $L$}{the real numbers of Theorem~1}
\nt{$a$}{a natural number: \m{\#I=3^{a+1}} in Theorem~2 (not the integer $a(D)$ of the data)}
\nt{$\mathcal S_1$, $\mathcal S_2$}{sets of polynomials with at most one element}
\nt{$\mathbb C[G]$, $y(g)$}{the complex group algebra of a finite group~$G$; the coefficient of~$y$ at~$g$}
\nt{$T$, $\theta_x$, $x_\chi$}{an injective character of $\bar{\F}^\times$; the character named by~$x$; the unit that names~$\chi$}
\nt{$\chi$}{a multiplicative character of~$\F$}
\nt{$\tau_P$, $\Xi_c$}{Green's class functions (Definition~\ref{d:idem})}
\nt{$F*H$, $\mathbb 1$, $\beta_x$, $W$, $V$}{the induction product, the constant function~$1$, the sum of $\theta_x$ over eigenvalues; a stable subspace of a space}
\nt{\m{\labelIdem{I}{c}}}{the idempotent of the label~$c$ in \m{\mathbb C[\mathrm{GL}_I(\F)]} (Definition~\ref{d:idem})}
\nt{\m{\tupleIdem{I}{\mathbf c}}}{the tensor product of the \m{\labelIdem{I}{c_v}} (Definition~\ref{d:tuple})}
\nt{\m{\goodTuples{M}{r}}}{the set of good tuples (Definition~\ref{d:good})}
\nt{$\mathcal P_v$, $\mathcal A_v$, $\mathcal B_v$, $\mathcal C$}{sets of monic polynomials in the proof of Theorem~1; a finite set of families}
\nt{$\varepsilon_0$}{the number \m{(1-r/q)\,q^{-r}}}
\nt{$\xi$, $\xi_0$}{roots of a polynomial in $\bar{\F}$}
\nt{$U_3$}{the group of cube roots of unity in~$\F$}
\nt{$\varphi$}{Euler's function}
\nt{$\pi$, $i$ in \m{2\pi i}}{the circle constant and the imaginary unit}
\nt{$n$, $l$, $o$, $p$, $d$, $i$, $s$, $z$, $\delta$}{auxiliary integers: exponents, orders, a prime, degrees, indices}
\nt{$g$, $x$, $y$}{a matrix or group element; an element of a field; an element of a group algebra}
\nt{$D$, $g_0(D)$, $g_1(D)$, $M_\omega(D)$, $b(D)$}{the data of the proof and objects attached to it, named in §\ref{s:use} only: two primes, an integer, a scale}
\end{multicols}
\endgroup

\head{References}
{\fontsize{9.4}{12.6}\selectfont
J.~A.~Green, \emph{The characters of the finite general linear groups}, Trans. Amer. Math. Soc. 80 (1955).
Named by the library's files as the source of Inputs 1 and~2. Place for Input~1, as located by the
second reading of this write-up: §8, Theorems 13--14, pp.~441--444. I did not open the paper.\par}

\end{document}
